\documentclass{ximera}

\title{Area Between Curves Activity Solution Guide}
\author{MATH 426: Calculus II}

\begin{document}
\begin{abstract}
   Working with the peers in your group, solve the following problems. Make sure to show and justify all your work. Make sure everyone in the group understands the solution and participates. Be prepared to report your answers to the whole class. 
\end{abstract}
\maketitle

\begin{exercise}
Consider the definite integral:
\begin{align}
	\int_{a}^b f(x) dx \nonumber
\end{align}
    \begin{enumerate}
        \item Assume that $a$ and $b$ are both constants with $a < b$. Assume $f(x)>0$. Draw a generic sketch of the meaning of this definite integral. \\
        \textcolor{blue}{Here is an example (note that to graph in Desmos, it cannot be generic, so a similar sketch without the numerical/function is what we are expecting):}
        \begin{center}
            \desmos{yt1y0w6asx}{800}{600}
        \end{center}
        \textcolor{blue}{Make sure that expression 1 is toggled ON and expressions 2 and 3 are toggled OFF (click the circle next to the expression).}
         Consider the definite integral:
\begin{align}
	\int_{a}^b g(x) dx \nonumber
\end{align}
        \item Assume that $a$ and $b$ are both constants with $a < b$. Assume that $g(x)$ has exactly one {\bf{root}} in the interval $(a,b)$. Draw a generic sketch of the meaning of this definite integral.\\
        \textcolor{blue}{Here is an example (note that to graph in Desmos, it cannot be generic, so a similar sketch without the numerical/function is what we are expecting):}
        \begin{center}
            \desmos{yt1y0w6asx}{800}{600}
        \end{center}
        \textcolor{blue}{Make sure that expression 1 is toggled OFF and expressions 2 and 3 are toggled ON (click the circle next to the expression).}

        \item Suppose that you need to compute the total {\bf{geometric}} area between $g(x)$ and the $x$-axis between $a$ and $b$. How would you construct this area using integrals? (Can you construct the area without using an absolute value?)\\
        \textcolor{blue}{The geometric area is always positive, so we \textit{could} use an absolute value like this: $\int_a^b |g(x)|dx$, but it will often help us to not use absolute value and instead break the integral up into pieces.  Assuming the root of $g(x)$ is at $x=c$ with $a<c<b$ and $g$ is first above, then below the $x$-axis, as in my graph:
        \[\int_a^c g(x)dx+\int_c^b -g(x)dx=\int_a^c g(x)dx-\int_c^b g(x)dx\]}
    \end{enumerate}

\end{exercise}





\begin{exercise}
    Have your group create two affine functions with positive slopes:
    \begin{align}
    	L_{1}(x) = m_{1}x + b_1, \hspace{.2in} L_{2}(x) = m_{2} x + b_2 \nonumber
    \end{align}
    Choose your slopes so that {\bf{both}} $m_{1} >m_2$, and $b_1>b_2$. Plot both of these lines on the interval $x\in [0,2]$. Your sketch should reveal that one of the lines is always on `top.'
    \begin{enumerate}
        \item Express the area under $L_{1}$ as a definite integral. Express the area under $L_2$ as a definite integral. Determine how to combine these integrals to capture the geometric area between them.\\
        \textcolor{blue}{Area under $L_1$: $\int_0^2 m_1 x+b_1dx$\\
        Area under $L_2$: $\int_0^2 m_2 x+b_2dx$\\
        Geometric area betwee $L_1$ and $L_2$: $\int_0^2 m_1 x+b_1dx-\int_0^2 m_2x+b_2dx$, or, equivalently, $\int_0^2 m_1x+b_1-(m_2x+b_2)dx$.}
        \item Have your group discuss the situation where you are trying to determine the geometric area between two curves $f(x)$ and $g(x)$ (not the same as your generic line sketches above) on a specific interval $(a,b)$. If you are struggling to identify what is going on, sketch the areas represented and relate them to integrals.
        \begin{enumerate}
            \item What happens if one of the functions has all positive values, and the other function has all negative values? \\
            \textcolor{blue}{The integral of the all postive function will be positive, and the integral of the all negative function will be negative.  We want to add the absolute values together, but this is the same as positive integral - (negative integral).}
            \item What happens if {\bf{both}} of the functions have negative values?\\
            \textcolor{blue}{If both functions are negative, then the integrals for both functions will be negative.  We want the positive area of the space between the two functions, so if we take integral of the function closer to the axis (negative and smaller) - integral of the function further from the axis (negative and larger) will result in the positive area between. }
            \item What happens if the functions cross over one another in the domain?\\
            \textcolor{blue}{If the functions cross, we will need to split into two integrals so we can always take top-bottom.}
            \item Discuss whether there is one framework which can handle all of the situations you considered separately.\\
            \textcolor{blue}{We can always use $\int_a^b |f(x)-g(x)|dx$, BUT this is not always straightforward or convenient to actually calculate.  A better framework comes from idenitifying the ``top'' and ``bottom'' functions and calculating $\int_{a_i}^{b_i} (top-bottom) dx$ for each interval.}
        \end{enumerate}
    \end{enumerate}

\end{exercise}

\begin{exercise}
    Use your discussion and your framework to \textbf{set up} the geometric area between each of the following sets of curves. (Use the interval $x\in (-2,2)$ for all of these exercises.) %If you have fully grasped the important concepts, you will not need to use symmetry or absolute value to correctly capture the areas.)
    \begin{align}
        f_{1} &= x^2 -1, \hspace{.2in} &g_{1} = 1-x^2 \nonumber \\ 
        f_{2} &= x+2, \hspace{.2in} &g_{2} = x-2 \nonumber\\
        f_{3} &= x, \hspace{.2in} &g_{3}= -x \nonumber\\
        f_{4} &= 2x, \hspace{.2in} &g_{4}= -\frac{x}{2}\nonumber\\
        f_5 &= x^2-2x+2, \hspace{.2in} &g_5 = x \nonumber
    \end{align}
    (Your TA/LA can provide numerical answers for you to check your constructions).\\
    \textcolor{blue}{ Area between $f_1$ and $g_1$: $\int_{-2}^{-1}(x^2-1)-(1-x^2)dx + \int_{-1}^1 (1-x^2)-(x^2-1)dx+\int_1^2 (x^2-1)-(1-x^2)dx$\\
    Area between $f_2$ and $g_2$: $\int_{-2}^2 (x+2)-(x-2)dx$\\
    Area between $f_3$ and $g_3$: $\int_{-2}^0 (-x-x)dx+\int_0^2 x-(-x)dx$\\
    Area between $f_4$ and $g_4$: $\int_{-2}^0 (-\frac{x}{2}-2x)dx +\int_0^2 2x-(-\frac{x}{2})dx$\\
    Area between $f_5$ and $g_5$: $\int_{-2}^1 (x^2-2x+2-x)dx+\int_1^2 x-(x^2-2x+2)dx$ }
\end{exercise}



\begin{exercise}
    Suppose you want to find the geometric area between two polynomial functions on a given interval $(a,b)$
        \begin{align}
        	P(x) = c_{n}x^n + c_{n-1}x^{n-1} + \dots +c_1x+ c_{0}, \hspace{.2in} Q(x) = k_{n}x^n + k_{n-1}x^{n-1} + \dots k_1x+ k_{0}.\nonumber
        \end{align}
With your groups discuss both of the following questions:
        \begin{enumerate}
            \item Under what conditions would this question be 'easy'?(Think about mathematical properties and situational properties, what would make the math easy to set-up, what would make the math easy to do?)\\
            \textcolor{blue}{We will need to be able to find any intersection points between the polynomials, so the simpler the polynomial forms, the easier the problem will be.  Also, we need to find out which function is on top and on the bottom, so the simpler the polynomials are, the easier to sub in values to figure this out.  Additionally, the number of intersections will complicate setting up the integrals. }
            \item Under what conditions would you be able to find intersections between $P$ and $Q$ by hand?\\
            \textcolor{blue}{Basically, we need to be able to solve the equation $P(x)=Q(x)$.  Unless nice factoring, cancelling or other nice things happen, we can do this by hand up to quadratic equations.}
        \end{enumerate}
\end{exercise}






\end{document}